\scanpage{14}
\mspage{12}
en d'autres termes
\[
 \varepsilon_{K'_n}(\Theta_n)=(\varepsilon^e)^{\unsure{n|\cdot}}_{K_n}(\Theta_n)
\]
et par suite
\[
 \varepsilon_{K'_n}\mid K_n=(\varepsilon^e)_{K_n}.
\]

Passant à la limite sur \(n\), on trouve la

\uline{Proposition} 3.12. Si \(\mathfrak mV'=\mathfrak m'{}^e\), alors l'hom. \inserted{\unsure{can|can.}} (3.8) provenant de (3.5) s'identifie à la multiplication par \(e\).

En particulier, si \(e=1\), cet hom. est l'identité, \deleted{\unsure{un|et}} et à fortiori un isomorphisme. \inserted{Ceci s'applique en particulier si \(V'=\widehat V\).} En fait, on a mieux :

\uline{Prop.} 3.13. Supposons \(V'=\widehat V\), \deleted{et \(V\) excellent} \deleted{et \(V\) excellent (i.e. \(\widehat K\) séparable sur \(K\))} \deleted{alors} Alors les hom de suites exactes (3.5) et (3.7) sont des isomorphismes.

\deleted{\inserted{\unsure{disons.|déjà vu.}} \(\pi'\unsure{\mathrel{-}|\longrightarrow}\pi\) est \deleted{\unsure{tjrs|toujours}} surjectif, i.e. si \(K_1\) est une extension finie séparable de \(K\), alors \(K_1\otimes_K\widehat K\) est un corps. \deleted{\unsure{Il faut|Il en résulte}} En effet, la clôture normale \(V_1\) de \(V\) dans \(K_1\) est \uline{finie} \deleted{\unsure{sur \(V\)|\(\widehat V\)}} par suite \(V_1\otimes_V\widehat V=\widehat V_1\) est intègre (car \(V_1\) est un anneau \inserted{de val. discrète;} donc \(K_1\otimes\widehat K=\unsure{\widehat V_1|V_1}\otimes\widehat K\) est un corps.}

\deleted{Comme \(\pi'_0\longrightarrow\pi_0\) est évidemment un isom, et aussi \(I'_t\mathrel{-}I_t\) \deleted{\unsure{par l'égalité|par l'équalité}} on est ramené en fait à examiner \(P'\mathrel{-}P\). Ceci \inserted{(ou directement \(\pi'\mathrel{-}\pi\), \deleted{\unsure{\(I\)|et}} donc \(I'\mathrel{-}I\))} est tjrs injectif, d'après le résultat connu de \deleted{\unsure{approximation|approximations}} \inserted{continuité} des solutions d'une équation algébrique p.r. aux coefficients, qui implique que des équations séparables \uline{voisines} définissent des extensions isomorphes de \(\widehat K\). \margin{\(\widetilde V\subset\widetilde{\widehat V}\subset\widehat{\widetilde V}\)}
% slanted
\deleted{D'autre part, c'est surjectif} \inserted{Comme \(I'\longrightarrow I\) est surjectif} (d'après le résultat \deleted{\unsure{que|qui} \illegible{1} fait l'hypothèse d'excellence}}

\deleted{\uline{Remarque} 3.14. La prop. 3.13. montre que dans la plupart des questions \deleted{\unsure{concernant|correspondant}} du séminaire, on pourrait supposer \(V\) \uline{complet}.}

% archive sheet number: 13
